\section{Robustness and inference}\label{sec:robustness}

This section subjects the main results to a battery of specification, inference, and
design checks, and then to an out-of-sample replication. The overall message is that the
wage null is very robust, the compositional results are moderately robust, and the
employment-level result is fragile.

\subsection{Specification robustness}

Table~\ref{tab:robspec} reports the $\mathrm{Low}\times\mathrm{Post}$ coefficient for the
four headline outcomes under eight alternative specifications: dropping the controls,
dropping the pandemic quarter of early 2021, adding department-specific linear trends,
restricting to prime-age workers, using an alternative education cutoff that removes the
categories at the boundary between the two groups, adding one-digit occupation and
industry fixed effects, and clustering by department and skill. The wage estimate remains
a tight null across every specification. The formality and self-employment estimates are
stable in sign and magnitude throughout, moving within a narrow band. The employment
estimate, by contrast, attenuates noticeably when the pandemic quarter is dropped or when
department-specific trends are included, which reinforces our reluctance to interpret it
causally.

\begin{table}[!tbp]\centering
\caption{Specification robustness of the DiD estimates}\label{tab:robspec}
\begin{threeparttable}
\footnotesize
\setlength{\tabcolsep}{4pt}
\begin{tabular}{lcccc}
\toprule
Specification & Log hourly wage & Employment & Formal empl. & Self-employment \\
\midrule
Baseline & -0.021 & -0.025$^{***}$ & -0.046$^{***}$ & 0.036$^{***}$ \\
 & (0.012) & (0.006) & (0.007) & (0.006) \\
 & \footnotesize[69,534] & \footnotesize[278,064] & \footnotesize[181,317] & \footnotesize[181,317] \\
No controls & -0.008 & -0.021$^{***}$ & -0.042$^{***}$ & 0.042$^{***}$ \\
 & (0.012) & (0.006) & (0.007) & (0.006) \\
 & \footnotesize[69,534] & \footnotesize[278,064] & \footnotesize[181,317] & \footnotesize[181,317] \\
Drop 2021Q1 (pandemic) & -0.023$^{*}$ & -0.014$^{***}$ & -0.042$^{***}$ & 0.031$^{***}$ \\
 & (0.012) & (0.005) & (0.008) & (0.005) \\
 & \footnotesize[65,279] & \footnotesize[259,135] & \footnotesize[168,908] & \footnotesize[168,908] \\
Department linear trends & -0.022$^{*}$ & -0.017$^{***}$ & -0.040$^{***}$ & 0.034$^{***}$ \\
 & (0.012) & (0.006) & (0.007) & (0.006) \\
 & \footnotesize[69,534] & \footnotesize[278,064] & \footnotesize[181,317] & \footnotesize[181,317] \\
Prime age (25--55) & -0.030$^{*}$ & -0.023$^{***}$ & -0.042$^{***}$ & 0.029$^{***}$ \\
 & (0.015) & (0.004) & (0.008) & (0.007) \\
 & \footnotesize[46,021] & \footnotesize[161,115] & \footnotesize[119,540] & \footnotesize[119,540] \\
Alternative education cutoff & -0.015 & -0.039$^{***}$ & -0.065$^{***}$ & 0.050$^{***}$ \\
 & (0.012) & (0.007) & (0.006) & (0.007) \\
 & \footnotesize[40,526] & \footnotesize[186,263] & \footnotesize[117,174] & \footnotesize[117,174] \\
Occupation \& industry FE & -0.007 & -0.016$^{***}$ & -0.036$^{***}$ & 0.027$^{***}$ \\
 & (0.011) & (0.003) & (0.006) & (0.005) \\
 & \footnotesize[69,534] & \footnotesize[213,810] & \footnotesize[181,317] & \footnotesize[181,317] \\
Cluster dept.$\times$skill & -0.021 & -0.025 & -0.046$^{***}$ & 0.036$^{***}$ \\
 & (0.013) & (0.020) & (0.007) & (0.007) \\
 & \footnotesize[69,534] & \footnotesize[278,064] & \footnotesize[181,317] & \footnotesize[181,317] \\
Low-skilled group linear trend & -0.021 & -0.005 & -0.004 & 0.006 \\
 & (0.019) & (0.008) & (0.010) & (0.009) \\
 & \footnotesize[69,534] & \footnotesize[278,064] & \footnotesize[181,317] & \footnotesize[181,317] \\
Include transition quarter (2022Q2) & -0.018 & -0.025$^{***}$ & -0.044$^{***}$ & 0.035$^{***}$ \\
 & (0.012) & (0.005) & (0.007) & (0.006) \\
 & \footnotesize[74,310] & \footnotesize[296,954] & \footnotesize[193,762] & \footnotesize[193,762] \\
Exclude agriculture (agrarian regime) & -0.022$^{*}$ & -- & -0.038$^{***}$ & 0.028$^{***}$ \\
 & (0.012) &  & (0.007) & (0.005) \\
 & \footnotesize[53,755] & \footnotesize[--] & \footnotesize[114,212] & \footnotesize[114,212] \\
Include public-sector workers & -0.002 & -- & -0.037$^{***}$ & 0.030$^{***}$ \\
 & (0.010) &  & (0.008) & (0.006) \\
 & \footnotesize[87,033] & \footnotesize[--] & \footnotesize[200,383] & \footnotesize[200,383] \\
\bottomrule
\end{tabular}
\begin{tablenotes}\footnotesize
\item Notes: Each cell is the $\mathrm{Low}\times\mathrm{Post}$ DiD coefficient (standard error clustered by department below; sample size in brackets) for the outcome in the column heading, under the specification in the row; all estimates are weighted by ENAHO survey weights, and Log hourly wage denotes the log real hourly wage. The baseline is equation~\eqref{eq:twfe}, which excludes the partially treated 2022Q2. The alternative education cutoff defines low-skilled as at most incomplete secondary and high-skilled as at least complete non-university tertiary, dropping the boundary categories. The agriculture exclusion and the public-sector inclusion apply to job-conditional outcomes only (industry and sector are undefined for non-workers). $^{*}p<0.1$, $^{**}p<0.05$, $^{***}p<0.01$.
\end{tablenotes}
\end{threeparttable}\end{table}


\subsection{Inference}

Because the treatment is assigned at the level of the skill group, we assess whether the
significance of the compositional results survives more demanding inference.
Table~\ref{tab:inference} collects the evidence. Columns (1) through (4) report
$p$-values for the baseline effect under department clustering, department-by-skill
clustering, the CR2 small-sample correction of \citet{pustejovsky_tipton_2018} applied to
department-level collapsed cells, and randomization inference that permutes the regional
bite of the minimum wage across departments. The self-employment effect is the most
resilient: it remains significant under department-by-skill clustering and under
randomization inference. The formality effect is significant under both clustering schemes
but only marginal under the most conservative CR2 and randomization procedures, which are
low-powered with twenty-five departments. The employment effect loses significance under
every procedure beyond simple department clustering. The wage null is, of course,
unaffected.

\begin{table}[!tbp]\centering
\caption{Inference robustness, design robustness, and out-of-sample validation}\label{tab:inference}
\begin{threeparttable}
\footnotesize
\begin{tabular}{lcccccc}
\toprule
Outcome & (1) Dept. & (2) Dept.$\times$skill & (3) CR2 & (4) RI (continuous) & (5) Continuous & (6) 2025 reform \\
\midrule
Log hourly wage & 0.106 & 0.126 & 0.854 & 0.029 & 0.017 & 0.002 \\
 & & & & & (0.012) & (0.011) \\
Employment & 0.000 & 0.210 & 0.507 & 0.003 & -0.034$^{***}$ & -0.002 \\
 & & & & & (0.012) & (0.005) \\
Formal empl. & 0.000 & 0.000 & 0.244 & 0.000 & -0.018$^{***}$ & -0.001 \\
 & & & & & (0.003) & (0.008) \\
Self-employment & 0.000 & 0.000 & 0.093 & 0.024 & 0.010$^{**}$ & -0.014$^{**}$ \\
 & & & & & (0.004) & (0.006) \\
\bottomrule
\end{tabular}
\begin{tablenotes}\footnotesize
\item Notes: Columns (1)--(3) report $p$-values for the baseline $\mathrm{Low}\times\mathrm{Post}$ effect under, respectively, department clustering (25 clusters), department$\times$skill clustering (50 clusters), and the CR2 small-sample correction \citep{pustejovsky_tipton_2018} applied to department-level collapsed cells. Column (4) reports the randomization-inference $p$-value for the CONTINUOUS-exposure design of column (5), not for the two-group $\mathrm{Low}\times\mathrm{Post}$ contrast, obtained by permuting the 25 regional Kaitz indices across departments (2{,}000 draws). Column (5) reports the continuous-exposure estimate, the coefficient on $\mathrm{Post}\times$(standardized department Kaitz index), with department-clustered standard errors in parentheses. Column (6) reports the $\mathrm{Low}\times\mathrm{Post}$ DiD estimate for the January-2025 reform (window 2024--2025), with department-clustered standard errors in parentheses. Log hourly wage denotes the log real hourly wage; all estimates are weighted by ENAHO survey weights. $^{*}p<0.1$, $^{**}p<0.05$, $^{***}p<0.01$.
\end{tablenotes}
\end{threeparttable}\end{table}


\subsection{Design robustness: bite intensity and triple differences}

Our identification rests on differential exposure by education, but the minimum wage also
bites differentially across regions, and this provides an independent source of variation.
Column (5) of Table~\ref{tab:inference} reports a continuous-exposure specification in the
spirit of \citet{card_1992}, which replaces the binary treatment with the standardized
department Kaitz index interacted with the post indicator. The results corroborate the
skill-based design: in departments where the minimum wage bites harder, formal employment
falls and self-employment rises by significantly more after the reform, while wages do not
respond. A triple-difference specification that interacts the skill contrast with an
indicator for high-bite departments is not statistically significant, which we read as a
lack of power in the binary regional split rather than as evidence against the mechanism,
given that the continuous version is informative.

\subsection{Sensitivity to violations of parallel trends}

Even flat pre-trends are estimated with error, so we apply the sensitivity analysis of
\citet{rambachan_roth_2023}. Figure~\ref{fig:honest} plots the robust confidence set for
the average post-reform effect on employment and formality as we allow the post-reform
violation of parallel trends to be as large as a multiple $\bar{M}$ of the largest
pre-reform violation. The formality effect remains bounded away from zero only for small
departures, up to about $\bar{M}=0.1$, and the self-employment effect only under near-exact
parallel trends, whereas the employment effect includes zero even under the smallest departure.
The compositional results are thus robust to conventional inference but sensitive to modest
trend violations, which reinforces our decision to emphasize the wage null, the one result
that is both precisely estimated and replicated out of sample.

\begin{figure}[H]
\centering
\includegraphics[width=0.8\textwidth]{fig7_honestdid.pdf}
\caption{Sensitivity to violations of parallel trends
\citep{rambachan_roth_2023}. Robust confidence sets for the average post-reform effect
under the relative-magnitudes restriction, which bounds the post-reform trend violation by
$\bar{M}$ times the largest pre-reform violation.}
\label{fig:honest}
\end{figure}

\subsection{Placebo reforms and leave-one-out}

We assign fake reforms to each interior quarter of the pre-reform window and re-estimate
the DiD using only pre-reform data. The wage placebos are uniformly null, the formality and
self-employment placebos are mostly null with one marginal case, but the employment placebos
are significant, which is a direct symptom of the pandemic-driven pre-trend in that outcome.
Leave-one-department-out and leave-one-quarter-out re-estimation leave the point estimates
essentially unchanged for all four outcomes, indicating that no single region or quarter
drives the results. These checks are reported in the appendix.

\subsection{Out-of-sample validation: the 2025 reform}

Finally, we replicate the entire design around the second reform, which raised the minimum
wage to 1,130 soles on 1 January 2025, using the 2024 and 2025 quarters. This reform is of
almost identical magnitude to the 2022 one but occurs in a calmer, post-pandemic
environment. Column (6) of Table~\ref{tab:inference} reports the estimates. The wage effect
is again a null, which strengthens our central finding considerably: across two independent
reforms in different macroeconomic conditions, the minimum wage did not raise the relative
wages of low-skilled workers. The compositional effects, however, do not replicate: the
formality and employment effects are null around the 2025 reform, and the self-employment
effect reverses sign. We take this seriously. It implies that the reallocation we document
around 2022 is specific to that episode and may reflect the unusual dynamics of the
post-pandemic labor market rather than a stable structural response to the minimum wage. The
one conclusion that survives out of sample is the wage null.
